\chapter{Keyboard}
\label{ch:keyboard}

\index{keyboard}

While the mouse is acme's primary interaction tool, the keyboard
handles text entry, cursor movement, and a small set of editing
shortcuts.  This standalone version of acme remaps the cursor keys,
Home, End, and Page keys to work in a standard, familiar way.

\section{Cursor Movement}
\label{sec:cursor}

\index{cursor movement}

\begin{longtable}{@{}lp{3.2in}@{}}
\toprule
Key & Action \\
\midrule
\endhead
\key{Left Arrow} & Move cursor one character left. \\
\key{Right Arrow} & Move cursor one character right. \\
\key{Up Arrow} & Move cursor up one line. \\
\key{Down Arrow} & Move cursor down one line. \\
\key{Ctrl+Left} & Move to beginning of line (same as \key{Ctrl+A}). \\
\key{Ctrl+Right} & Move to end of line (same as \key{Ctrl+E}). \\
\key{Ctrl+A} & Move to beginning of line. \\
\key{Ctrl+E} & Move to end of line. \\
\bottomrule
\end{longtable}

\subsection{Sticky Column}
\index{sticky column}

The up and down arrow keys maintain a \emph{sticky column} position.
When you press \key{Up} or \key{Down}, acme remembers the horizontal
pixel position of the cursor and tries to place it at the same
horizontal position on the adjacent line.  If the adjacent line is
shorter, the cursor goes to the end of that line, but the original
column is remembered.  When you reach a longer line again, the cursor
returns to the remembered column.

The sticky column resets whenever you press any key other than \key{Up}
or \key{Down}, or when you click the mouse.

\subsection{Scrolling with Arrow Keys}
\index{scrolling!arrow keys}

If the cursor is on the \emph{last visible line} and you press
\key{Down}, acme scrolls down one line and places the cursor on the new
last visible line.  Similarly, pressing \key{Up} on the first visible
line scrolls up one line.  This gives smooth, continuous scrolling while
navigating with arrow keys.

\section{Page Navigation}
\label{sec:paging}

\index{scrolling!page}

\begin{longtable}{@{}lp{3.2in}@{}}
\toprule
Key & Action \\
\midrule
\endhead
\key{Page Up} & Scroll up one full screen. \\
\key{Page Down} & Scroll down one full screen. \\
\key{Ctrl+Up} & Scroll up one full screen (same as Page~Up). \\
\key{Ctrl+Down} & Scroll down one full screen (same as Page~Down). \\
\bottomrule
\end{longtable}

\key{Page Down} does nothing if the end of the file is already visible.

\section{Jump to Position}
\label{sec:jump}

\index{Home key}
\index{End key}

\begin{longtable}{@{}lp{3.2in}@{}}
\toprule
Key & Action \\
\midrule
\endhead
\key{Home} & Move cursor to first character visible on screen. \\
\key{End} & Move cursor to beginning of last visible line. \\
\key{Ctrl+Home} & Jump to beginning of file. \\
\key{Ctrl+End} & Jump to end of file. \\
\key{Ctrl+Page Up} & Jump to beginning of file (same as \key{Ctrl+Home}). \\
\key{Ctrl+Page Down} & Jump to end of file (same as \key{Ctrl+End}). \\
\bottomrule
\end{longtable}

\key{Home} and \key{End} do \emph{not} scroll the text---they move the
cursor within the currently visible portion.  \key{Ctrl+Home},
\key{Ctrl+End}, \key{Ctrl+Page Up}, and \key{Ctrl+Page Down} scroll
as needed to show the target position.

\section{Editing Keys}
\label{sec:editkeys}

\index{Backspace}
\index{Delete key}

\begin{longtable}{@{}lp{3.2in}@{}}
\toprule
Key & Action \\
\midrule
\endhead
\key{Backspace} & Delete the character to the left of the cursor. \\
\key{Delete} & Delete the character to the right of the cursor. \\
\key{Ctrl+H} & Delete character left (same as \key{Backspace}). \\
\key{Ctrl+W} & Delete word to the left of the cursor. \\
\key{Ctrl+U} & Delete from cursor to beginning of line. \\
\bottomrule
\end{longtable}

\key{Backspace} deletes the character to the left, which is the
standard behavior on modern systems.  \key{Delete} (forward delete)
removes the character to the right.

\section{Text Entry}

\index{auto-indent}
\index{Tab key}

Typing inserts text at the cursor position, replacing any selection.
This is standard behavior---there are no modes.

\paragraph{Auto-indent.}
When auto-indent is enabled (via the \cmd{-a} flag or the \cmd{Indent}
command), pressing \key{Enter} inserts a newline followed by the same
leading whitespace as the previous line.

\paragraph{Tab.}
The \key{Tab} key inserts a literal tab character.  The display width of
a tab is controlled by the \cmd{Tab} command or the \env{tabstop}
environment variable (default: 4 columns).

\section{Unicode Composition}
\label{sec:compose}

\index{Unicode composition}
\index{Alt key}

Acme supports entering Unicode characters through a multi-key
composition system.  To compose a character:

\begin{enumerate}
\item Press and release the \key{Alt} key (this toggles composition mode).
\item Type the composition sequence (one or more characters).
\item The composed character is inserted.
\end{enumerate}

If the sequence is unrecognized, the literal characters are inserted.
Press \key{Alt} again to cancel composition mode.

\subsection{Common Compositions}

\begin{longtable}{@{}llll@{}}
\toprule
Sequence & Result & & Description \\
\midrule
\endhead
\texttt{'e} & \'{e} & & Acute accent \\
\texttt{`e} & \`{e} & & Grave accent \\
\texttt{\^{}e} & \^{e} & & Circumflex \\
\texttt{"e} & \"{e} & & Diaeresis / umlaut \\
\texttt{\~{}n} & \~{n} & & Tilde \\
\texttt{,c} & \c{c} & & Cedilla \\
\texttt{oe} & \oe & & Ligature \\
\texttt{*a} & $\alpha$ & & Greek alpha \\
\texttt{*b} & $\beta$ & & Greek beta \\
\texttt{*p} & $\pi$ & & Greek pi \\
\texttt{!=} & $\neq$ & & Not equal \\
\texttt{<=} & $\leq$ & & Less or equal \\
\texttt{>=} & $\geq$ & & Greater or equal \\
\texttt{+-} & $\pm$ & & Plus-minus \\
\texttt{->} & $\rightarrow$ & & Right arrow \\
\texttt{<-} & $\leftarrow$ & & Left arrow \\
\texttt{sr} & $\sqrt{}$ & & Square root \\
\texttt{if} & $\infty$ & & Infinity \\
\bottomrule
\end{longtable}

\noindent
The full composition table contains over 500 entries, covering Latin
accented characters, Greek, Cyrillic, mathematical symbols, currency
symbols, arrows, and more.  See Appendix~\ref{app:unicode} for the
complete table.

\section{Summary}

\begin{longtable}{@{}llp{2.5in}@{}}
\toprule
Category & Keys & Notes \\
\midrule
\endhead
Movement & Arrows, \key{Ctrl}+Arrows & Sticky column on up/down \\
Pages & \key{PgUp}, \key{PgDn}, \key{Ctrl}+Up/Dn & Full-screen scroll \\
Jump & \key{Home}, \key{End} & Within visible area \\
File jump & \key{Ctrl+Home}, \key{Ctrl+End} & Beginning / end of file \\
Delete & \key{Backspace}, \key{Delete} & Left / right of cursor \\
Words & \key{Ctrl+W} & Delete word left \\
Lines & \key{Ctrl+U}, \key{Ctrl+A}, \key{Ctrl+E} & Delete/move to line boundary \\
Compose & \key{Alt}, then sequence & Unicode character entry \\
\bottomrule
\end{longtable}
